\documentclass[10pt]{article}
\providecommand{\FIGDIR}{fig}
\newcommand{\DIGESTDATE}{31 August 2026}
% ===================== PM Research Watch - shared preamble =====================
\usepackage[T1]{fontenc}
\usepackage[utf8]{inputenc}
\usepackage{amsmath,amssymb}
\usepackage{tgpagella}
\usepackage{tgheros}
\renewcommand{\sfdefault}{qhv}
\renewcommand{\ttdefault}{lmtt}
\usepackage{microtype}
\usepackage[english]{babel}

\usepackage[a4paper,top=17mm,bottom=17mm,left=17mm,right=17mm,
            headsep=6mm,footskip=10mm]{geometry}
\usepackage{graphicx}
\usepackage{xcolor}
\usepackage{colortbl}
\usepackage{tikz}
\usetikzlibrary{calc}
\usepackage[most]{tcolorbox}
\usepackage{titlesec}
\usepackage{fancyhdr}
\usepackage{booktabs}
\usepackage{array}
\usepackage{longtable}
\usepackage{enumitem}
\usepackage{multicol}
\usepackage{ragged2e}
\usepackage[hidelinks]{hyperref}

% ------------------------------------------------------------------ palette
\definecolor{Ink}    {HTML}{14262E}
\definecolor{Deep}   {HTML}{0F4C5C}
\definecolor{Teal}   {HTML}{1F7A8C}
\definecolor{Sky}    {HTML}{4FA3C4}
\definecolor{Sage}   {HTML}{6FA287}
\definecolor{Amber}  {HTML}{E3A008}
\definecolor{Clay}   {HTML}{B5643C}
\definecolor{Coral}  {HTML}{D1495B}
\definecolor{Violet} {HTML}{6A5B9E}
\definecolor{Slate}  {HTML}{4A6572}
\definecolor{Paper}  {HTML}{FBF9F5}
\definecolor{Mist}   {HTML}{EEEAE2}
\definecolor{Rule}   {HTML}{D9D3C9}

\pagecolor{Paper}
\color{Ink}

\hypersetup{colorlinks=true, linkcolor=Deep, urlcolor=Teal, citecolor=Teal,
            pdftitle={PM Research Watch}, pdfauthor={Automated daily digest}}

% ------------------------------------------------------------------ headings
\titleformat{\section}
  {\sffamily\bfseries\large\color{Deep}}{}{0pt}
  {}[{\vspace{1.2mm}\color{Amber}\titlerule[1.1pt]}]
\titlespacing*{\section}{0pt}{5.5mm}{2.8mm}

\titleformat{\subsection}
  {\sffamily\bfseries\normalsize\color{Teal}}{}{0pt}{}
\titlespacing*{\subsection}{0pt}{3.6mm}{1.4mm}

\setlist[itemize]{leftmargin=4.2mm, itemsep=1.1mm, topsep=1.3mm, parsep=0pt,
                  label=\textcolor{Teal}{\footnotesize$\blacktriangleright$}}

% ------------------------------------------------------------------ header/footer
\pagestyle{fancy}
\fancyhf{}
\renewcommand{\headrulewidth}{0pt}
\renewcommand{\footrulewidth}{0pt}
\fancyhead[L]{\sffamily\scriptsize\color{Slate}PM RESEARCH WATCH\;\textcolor{Rule}{\textbar}\;\DIGESTDATE}
\fancyhead[R]{\sffamily\scriptsize\color{Slate}Particulate matter monitoring \& health effects}
\fancyfoot[C]{\sffamily\scriptsize\color{Slate}\thepage}

% ------------------------------------------------------------------ components
% metric tile
\newcommand{\metric}[3]{%
  \begin{tcolorbox}[enhanced, colback=white, colframe=Rule, boxrule=0.5pt,
    arc=1.6mm, left=2.2mm, right=2.2mm, top=1.8mm, bottom=1.8mm,
    borderline west={1.7mm}{0pt}{#3}, nobeforeafter, valign=center,
    height=17mm, width=\linewidth]
    \RaggedRight
    {\sffamily\bfseries\fontsize{17}{17}\selectfont\color{#3}#1}\\[0.5mm]
    {\sffamily\fontsize{6.5}{7.4}\selectfont\color{Slate}#2}
  \end{tcolorbox}}

% pull-quote / key insight box
\newtcolorbox{keybox}[1][]{enhanced, colback=Mist, colframe=Mist,
  arc=1.6mm, left=3.2mm, right=3.2mm, top=2.4mm, bottom=2.4mm,
  borderline west={1.7mm}{0pt}{Amber}, #1}

% subtopic band
\newcommand{\band}[2]{%
  \vspace{3.2mm}\noindent
  \begin{tcolorbox}[enhanced, colback=#2, colframe=#2, arc=1.2mm,
    left=2.6mm, right=2.6mm, top=1.1mm, bottom=1.1mm, boxrule=0pt,
    width=\linewidth, nobeforeafter]
    {\sffamily\bfseries\small\color{white}#1}
  \end{tcolorbox}\vspace{1.6mm}}

% one paper entry: {short cite}{title}{journal}{doi}{body}
\newcommand{\paperentry}[5]{%
  \par\vspace{1.9mm}\noindent
  {\sffamily\bfseries\footnotesize\color{Ink}#1}\;%
  {\sffamily\fontsize{7}{8}\selectfont\color{Slate}#3}\\[0.4mm]
  {\footnotesize\itshape\color{Deep}#2}\\[0.6mm]
  {\footnotesize #4}\\[0.3mm]
  {\fontsize{6.6}{7.6}\selectfont\color{Slate}\texttt{doi:}\,\href{https://doi.org/#5}{\texttt{\detokenize{#5}}}}
  \par}

% caption style
\newcommand{\figcap}[1]{\par\vspace{1.0mm}
  {\fontsize{7.2}{8.6}\selectfont\color{Slate}\textsf{\textbf{Fig.}}\; #1\par}\vspace{1.5mm}}

\providecommand{\permil}{\textperthousand}

\setlength{\parindent}{0pt}
\setlength{\parskip}{0pt}
\renewcommand{\arraystretch}{1.18}


\begin{document}

% ==================================================================== MASTHEAD
\thispagestyle{empty}
\begin{tikzpicture}[remember picture, overlay]
  \fill[Deep] (current page.north west) rectangle
      ($(current page.north east)+(0,-67mm)$);
  \fill[Coral] ($(current page.north west)+(0,-67mm)$) rectangle
      ($(current page.north east)+(0,-68.4mm)$);
  \foreach \x/\y/\r in {22/12/0.9, 41/26/1.5, 63/9/0.7, 88/31/1.1, 112/17/1.9,
                        138/29/0.8, 159/13/1.3, 178/54/1.0, 196/21/1.6,
                        30/44/0.6, 74/57/1.2, 125/48/0.9, 168/6/0.8, 191/59/0.7,
                        52/18/1.0, 100/60/0.6, 148/20/1.1, 66/49/1.4, 205/45/0.8}
    \fill[Sky, opacity=0.30] ($(current page.north west)+(\x mm,-\y mm)$) circle (\r mm);
\end{tikzpicture}

\vspace*{4mm}
{\sffamily\fontsize{7.6}{9}\selectfont\color{Coral}%
 MONTHLY ROLLUP\;\textbullet\;PhD RESEARCH WATCH\par}
\vspace{2.2mm}
{\sffamily\bfseries\fontsize{27}{29}\selectfont\color{white}%
 Particulate Matter\par}
{\sffamily\fontsize{27}{29}\selectfont\color{Sky}%
 Monitoring \& Health Effects\par}
\vspace{5.0mm}
{\sffamily\fontsize{8.4}{10}\selectfont\color{white}%
 August 2026\;\;\textcolor{Coral}{\textbar}\;\;Entry window 1--31 August 2026, 31 daily issues\;\;%
 \textcolor{Coral}{\textbar}\;\;556 records pooled\par}
\vspace{18mm}

% -------------------------------------------------------------- metric strip
\noindent
\begin{minipage}[t]{0.190\linewidth}\metric{556}{RECORDS POOLED\\(31 DAILY ISSUES)}{Deep}\end{minipage}\hfill
\begin{minipage}[t]{0.190\linewidth}\metric{140}{SENSING /\\INSTRUMENTATION (25\%)}{Teal}\end{minipage}\hfill
\begin{minipage}[t]{0.190\linewidth}\metric{148}{EFFECT ESTIMATES\\FROM 63 SOURCES}{Sky}\end{minipage}\hfill
\begin{minipage}[t]{0.190\linewidth}\metric{48}{RECORDS FROM\\AMT + ACP}{Amber}\end{minipage}\hfill
\begin{minipage}[t]{0.190\linewidth}\metric{12}{ULTRAFINE /\\NANOSCALE (2\%)}{Clay}\end{minipage}

\vspace{6mm}

% ------------------------------------------------------------ signal of month
\section*{\sffamily\bfseries\color{Deep}Signal of the month}
\vspace{-1.5mm}{\color{Coral}\rule{\linewidth}{1.1pt}}\vspace{2.5mm}

\begin{keybox}
{\sffamily\bfseries\small\color{Deep}Low-cost PM sensor calibration stopped being a
collection of one-site case studies in August. Eight independent groups published work
whose common subject is not \emph{whether} a correction works but \emph{how far it
travels} --- across cities, across seasons, across distance from the nearest
reference.}\\[1.4mm]
\footnotesize
The clearest statement is Ginsburg et al.\ (\textit{Aerosol Sci Technol}), who tested
whether the US EPA PurpleAir correction transfers worldwide and found it does not
without regional recalibration. Zhivkov et al.\ (\textit{Atmosphere}) put a number on the
spatial half of the same problem: using 24 months from Sofia with four co-located and
eighteen remote sensors, calibration uncertainty grows \textbf{3.84--5.62\% per kilometre}
from the nearest reference station, so co-location sites 6\,km apart hold the network
midpoint inside a 30\% uncertainty increase --- a design rule, derived empirically, that a
network planner can actually use. Wang et al.\ (\textit{Environ Sci Technol}) extended
correction to a nationwide dense network without local references at all; Balasubramanian
et al.\ (\textit{Environ Pollut}) built the US-wide PurpleAir model; Blaga et al.\
(\textit{PLOS Clim}) took a city-scale network to regulatory-compliant output; Chung
et al.\ (\textit{Sensors}) characterised PurpleAir \emph{longevity} and uncertainty across
2016--2025 deployments, which is the failure mode nobody calibrates for; Barkjohn et al.\
(\textit{Atmosphere}) shipped ASNAT, an R-Shiny QA/QC and reference-comparison tool, i.e.\
the tooling layer arriving after the methods; and G\"abel \& Hertig (\textit{AMT}) tied
recalibration practice to the current sensor-technology literature.
Two constraints published alongside them matter more than any of the corrections.
Chyzhykov et al.\ (\textit{J Hazard Mater}) quantified \textbf{particle-bound water as a
bias on gravimetric PM mass} --- that is the \emph{reference} being wrong, not the sensor
--- and Benjamin et al.\ (\textit{Sensors}) showed environmental extrapolation failing
outside the training envelope for both CO and PM$_{2.5}$ nodes. \textbf{The month's
verdict: transferability, drift and reference validity are now the three named open
problems, and the field has stopped pretending a single regression solves any of them.}
\end{keybox}

\vspace{3mm}

\section*{\sffamily\bfseries\color{Deep}What changed in August}
\vspace{-1.5mm}{\color{Coral}\rule{\linewidth}{1.1pt}}\vspace{2.5mm}

\footnotesize
\begin{itemize}
\item \textbf{The pipeline gap July's rollup named was closed, and it is the largest
single change in this corpus.} July's monthly stated that there was ``no date-windowed
channel into \textit{Atmos.\ Meas.\ Tech.} or \textit{Atmos.\ Chem.\ Phys.}'' and
proposed a Crossref by-journal feed keyed on \texttt{created}. That feed was built and
widened through the month. August's result: \textbf{Atmos Chem Phys is the single largest
venue in the corpus at 37 records, with Atmos Meas Tech at 11} --- 48 together, ahead of
\textit{Environ Sci Technol} (25) and \textit{Environ Pollut} (22). The pure
atmospheric-measurement literature that was structurally invisible in July is now the
backbone of the sensing cluster.

\item \textbf{Sensing is now the largest subtopic, at 140/556 (25\%).} July recorded 23.
Most of that difference is the AMT/ACP channel plus a full 31-issue cadence rather than
five issues, so it is \textbf{not} evidence that the field grew --- it is evidence of how
badly the July corpus under-sampled instrumentation. Exposure assessment \& modelling is
second at 129 (23\%). Together the two non-clinical clusters hold 48\% of the month.

\item \textbf{Ultrafine particles are still the metric nobody measures.} Twelve records
(2\%) use an ultrafine or nanoscale metric, \emph{down} from July's 5\% share. Yet Zhong
et al.\ (\textit{Environ Int}) showed photochemical ultrafine particles posing an
inflammatory risk \emph{at low PM$_{2.5}$ mass}; Rashid et al.\ (\textit{Toxics})
published a scoping review of UFP mortality evidence and exposure assessment; Blanco
et al.\ (\textit{JESEE}) showed mobile-monitoring design choices propagating into UFP
exposure models; and Zhang et al.\ (\textit{Environ Sci Technol}) built a wearable that
measures particle deposition in the human respiratory tract directly. The prediction
made in July's rollup --- that a field calibration of a low-cost particle-number or
black-carbon sensor against a CPC or aethalometer reference would appear within two
monthly issues --- \textbf{did not come true in August}. The nearest approaches are
Kirchhoff et al.'s CPC working-fluid evaluation (\textit{AMT}) and Shlipak et al.'s
source-apportioned black carbon in Addis Ababa, neither of which is a low-cost-instrument
validation. That prediction now carries into September, and one more miss should be read
as a hardware-cost bottleneck rather than a research-attention one.

\item \textbf{Effect estimates: 148 from 63 sources, and 110 of them exclude the null in
the harmful direction.} Only six were protective, and every one of those has a benign
explanation (an ozone term, a wildfire-season influenza-like-illness contrast, thyroxine
and menopause-timing shifts). Thirty-two cross the null. The largest ratios are all
extreme-contrast designs --- Xu et al.\ report OR $\approx$ 4.3--4.4 for cerebral
hypoperfusion against PM$_{2.5}$ and black carbon, and Li et al.\ OR 5.08 (3.45--7.47)
for paediatric allergic-rhinitis exacerbation comparing PM$_{2.5}$ $>$75 against
$<$35\,$\mu$g\,m$^{-3}$ --- and are not comparable to per-10\,$\mu$g\,m$^{-3}$ estimates.

\item \textbf{Compound and multi-stressor designs consolidated.} Mandalapu et al.\
(medRxiv) mapped joint heat and PM$_{2.5}$ exposure across US metropolitan areas against
historical HOLC redlining grades (\textbf{OR 3.97, 2.79--5.66} for hotspot residence in
the worst-graded areas); Yu et al.\ (\textit{Environ Sci Technol}) gave a probabilistic
treatment of PM$_{2.5}$--ozone co-occurring extremes under wildfire smoke; Wang L et al.\
closed the month with a lipid-fraction $\times$ PM$_{2.5}$ interaction on stroke. Single
-pollutant, single-endpoint time-series work continues to read as the low-information
tail.

\item \textbf{Sub-Saharan Africa quadrupled, from 4 records in July to 16, and it still
is not enough.} Shlipak et al.\ (Addis Ababa source apportionment), Mushore et al.\
(road-edge effect across three campaigns), Rehbein et al.\ (low-cost network source
triangulation in Ile-Ife) and Makacha et al.\ (personal PM$_{2.5}$ exposure determinants)
are all methodologically serious. Against 121 records from China, the imbalance remains
the corpus's most persistent structural feature.

\item \textbf{A dedup defect was found by this rollup and has been gated.} See the
provenance box: twelve records were shipped twice.
\end{itemize}

\clearpage

% ==================================================================== FIGURES
\section*{\sffamily\bfseries\color{Deep}Corpus analytics --- August 2026}
\vspace{-1.5mm}{\color{Coral}\rule{\linewidth}{1.1pt}}\vspace{2.5mm}

\noindent\includegraphics[width=\linewidth]{\FIGDIR/f1_subtopics.png}
\figcap{Subtopic distribution across all 556 pooled records. All ten subtopics are
populated, as they were in every week of the month. \textit{Sensing, forecasting \&
instrumentation} (140) and \textit{Exposure assessment \& modelling} (129) together take
48\% of the corpus --- a share that reflects the new AMT/ACP harvesting channel as much
as the underlying literature. \textit{Cardiovascular \& metabolic} (23) and
\textit{Reproductive \& developmental} (19) are the thinnest clusters. Assignment is
single-label by dominant contribution.}

\vspace{3mm}

\noindent\includegraphics[width=\linewidth]{\FIGDIR/w1_subtopic_trend.png}
\figcap{Records per issue date by subtopic across the 31 August issues. Unlike July's
equivalent panel this \emph{is} a genuine daily series --- every bar is one entry-date
window, with no backfill batch. Day-to-day variation is large (8 records on 29 August, 36
on 27 August) and is driven by deposit timing at a handful of high-volume publishers
rather than by anything in the literature. The ISO-week totals are steadier: 56, 127,
132, 111, 114 and 16 for the part-week ending 31 August. Read week totals, not day
totals.}

\vspace{3mm}
\noindent\includegraphics[width=\linewidth]{\FIGDIR/f2a_architecture.png}
\figcap{Study architecture over the month. Measurement campaigns lead at 113, ahead of
modelling/inventory (100), reviews (58) and observational cohorts (53). \textbf{The
``Metadata only'' group at 65 records (12\%) is the number to hold onto}: those are
records whose abstract was not deposited anywhere reachable at issue time. They are
counted in every total in this issue but carry no assessed content.
\textbf{Trials or interventions: 4 records in 556.} That ratio has not improved since
July and is the standing weakness of the evidence base.}

\clearpage

%% f2b_endpoint.png is DELIBERATELY OMITTED from this issue (2026-09-01).
%% At rollup scale the figure is 1087x2727 px: at \linewidth it requests 1256pt of
%% height against a ~700pt text block and overflowed page 6 by 507pt; height-capped
%% to 0.80\textheight it fits but renders ~90 category rows at roughly 3pt and is
%% unreadable. The root cause is the open ENDPOINT_CANON defect (logged 15 Aug):
%% `endpoint` is free text, so 556 records carry 144 distinct strings and the figure
%% has no bounded category set to plot. Shipping an illegible chart is worse than
%% shipping the numbers in prose, so the distribution is given as text below.
\begin{keybox}
{\sffamily\bfseries\footnotesize\color{Deep}Health endpoint distribution --- reported as
text, not as a figure}\\[1.2mm]
\footnotesize
The endpoint figure is \textbf{omitted from this issue on purpose}. The \texttt{endpoint}
field is free text and has never been canonicalised --- an open defect since 15 August ---
so the 556 pooled records carry \textbf{144 distinct endpoint strings}, and the pooled
chart degenerates into roughly ninety unreadable rows. The distribution itself is short
enough to state:
\textbf{210 records (38\%) carry no human health endpoint at all}, split across
``None (monitoring)'' 75, ``No health endpoint'' 71, ``None (exposure)'' 33 and
``None (instrument)'' 31 --- four spellings of one category, which is the defect in
miniature. A further \textbf{31 records have an endpoint that was not extracted}.
Of the remainder the named categories are \textbf{respiratory 36, cardiovascular 22,
other clinical 14, reproductive 13, neuro/cognitive 11, in-vitro toxicological 11,
mortality 10, aerosol chemistry 9, oncologic 9}, with the balance spread thinly across
more than a hundred one- and two-record strings. \textbf{Fix required before the
September rollup:} a bounded \texttt{ENDPOINT\_CANON} map in \texttt{plots.py}, in the
same pattern as \texttt{design\_group} and \texttt{geo\_group}, both of which already
work.
\end{keybox}

\vspace{3mm}

\noindent\includegraphics[width=\linewidth]{\FIGDIR/f3_metric_geography.png}
\figcap{Left --- particle metric. PM$_{2.5}$ alone in 213 records, unspeciated PM or
emissions in 160, PM$_{2.5}$ with PM$_{10}$ in 83. \textbf{Ultrafine or nanoscale metrics
appear in 12 records (2\%) and size-resolved number distributions in 11}, against 31 for
composition/speciation and 20 for optical properties. Right --- geography.
\textit{Global / multi-region} 126, China 121 (22\%), Europe 91, North America 82, East
Asia excluding China 32, South Asia 28, Southeast Asia 20, Middle East \& N.\ Africa 18,
\textbf{sub-Saharan Africa 16}, Latin America 13, Oceania 6, polar/remote marine 2,
Central Asia 1.}

\clearpage

\noindent\includegraphics[width=\linewidth]{\FIGDIR/f4_heatmap.png}
\figcap{Subtopic $\times$ study-architecture cross-tabulation over the month. The dense
cells are expected: sensing $\times$ measurement campaign, exposure $\times$
modelling/inventory. The informative absences are the same two the daily issues keep
reporting --- the trial column is nearly empty across every health subtopic, and the
sensing and exposure rows contribute almost nothing to the effect-estimate pool.}

\vspace{3mm}

\noindent\includegraphics[width=\linewidth]{\FIGDIR/f5_forest.png}
\figcap{Quantitative effect estimates extracted across August, panel 1 of 5.
\textbf{This is not a meta-analysis and no pooled summary is drawn.} Exposure contrasts
are not commensurable --- per-IQR, per-SD, per-10\,$\mu$g\,m$^{-3}$, threshold, quartile
and mixture-index contrasts all appear --- and the outcome scale mixes RR, OR, aOR, HR,
IRR, GMR, PR and time ratios across unrelated diseases. \textbf{Heterogeneity is
therefore not estimable and none is reported.} Of the 148 estimates, 110 exclude the null
in the harmful direction, 6 in the protective direction and 32 cross it.}

\clearpage

\noindent\includegraphics[width=\linewidth]{\FIGDIR/f5_forest_2.png}
\figcap{Effect estimates, panel 2 of 5. Same caveats as panel 1.}

\vspace{3mm}

\noindent\includegraphics[width=\linewidth]{\FIGDIR/f5_forest_3.png}
\figcap{Effect estimates, panel 3 of 5. Same caveats as panel 1.}

\clearpage

\noindent\includegraphics[width=\linewidth]{\FIGDIR/f5_forest_4.png}
\figcap{Effect estimates, panel 4 of 5. Same caveats as panel 1.}

\vspace{3mm}

\noindent\includegraphics[width=\linewidth]{\FIGDIR/f5_forest_5.png}
\figcap{Effect estimates, panel 5 of 5. Same caveats as panel 1. Read the panels as a
catalogue of what was reported, not as a ranking: the visually largest estimates are
uniformly extreme-contrast designs (threshold comparisons, redlining-grade contrasts,
mixture indices) and the tightest intervals belong to the largest cohorts.}

\clearpage

\noindent\includegraphics[width=\linewidth]{\FIGDIR/f6_lifecourse.png}
\figcap{Life-course windows addressed across the month. Counts are non-exclusive and
cover only records with an identifiable human life-course window; \textit{in vitro},
ex-vivo, animal and pure-instrumentation records are excluded, which is why the totals
are far below 556. Adolescence remains the thinnest window, as it was in July and in
almost every individual daily issue.}

\vspace{4mm}

\section*{\sffamily\bfseries\color{Deep}Emerging, fading, and recurring}
\vspace{-1.5mm}{\color{Coral}\rule{\linewidth}{1.1pt}}\vspace{2.5mm}

\footnotesize
\textbf{Emerging.} \textbf{(i) Calibration transferability as a named problem} --- see
the signal box; the shift is from ``here is a correction for our site'' to ``here is how
far a correction travels, and here is the uncertainty as a function of distance,
season and instrument age''. \textbf{(ii) The reference itself under scrutiny} ---
Chyzhykov et al.\ on particle-bound water biasing gravimetric mass, Anon et al.\ on a
gravimetric-method assessment campaign for PM$_{10}$ and PM$_{2.5}$, and Alexa et al.\ on
separating gaseous nitric acid from particulate nitrate. When the calibration target is
itself biased, every downstream $R^2$ inherits it. \textbf{(iii) Toxicity- and
source-weighted exposure} --- Chen et al.\ on source-specific oxidative potential of
PM$_{2.5}$ in Xi'an, Kang et al.\ on vehicle-origin condensable PM, and Song et al.'s
closing perspective on shifting control from mass to toxicity. \textbf{(iv) Citizen
science formalised} --- Marcon et al.\ (\textit{Curr Environ Health Rep}) reviewed
methodological approaches and effectiveness, and Cho et al.\ ran a community-scientist-led
mobile campaign; this had been anecdotal in July.

\vspace{2mm}
\textbf{Fading.} With 31 genuine daily windows this month, a decline claim is finally
defensible in principle --- but one month against one partly-backfilled month is still
not a trend, so nothing is asserted here as fading. The weakest-signal observation, and
it is only an observation: \textbf{single-city, single-pollutant time-series studies
continued to appear and continued to add nothing}, generally reporting PM as the weakest
term in their own multi-pollutant models. That was July's observation too, unchanged.

\vspace{2mm}
\textbf{Recurring groups and venues.} \textbf{No author-cluster analysis is possible in
this corpus and none is offered.} Records carry a first-author surname but no
affiliation field, so apparent repeats (Li, Zhang, Wang, Chen, Liu) are surname frequency
in a China-heavy corpus, not the same investigators. This has been an open defect since
15 August and remains one. Venue concentration is the analysis that \emph{is} supportable:
\textit{Atmos Chem Phys} 37, \textit{Atmos Environ} 34, \textit{Environ Sci Technol} 25,
\textit{Environ Pollut} 22, \textit{Environ Monit Assess} 17, \textit{Atmos Pollut Res}
14, \textit{Atmos Meas Tech} 11, \textit{Atmosphere} 11, \textit{J Hazard Mater} 10,
\textit{Environ Int} 8. Preprints contributed 31 records --- Research Square 21 under
\emph{three} spellings, medRxiv 6 under two, arXiv 4 --- which is a venue-normalisation
defect, logged below. Citation counts on records days to weeks old are uninformative and were not used
to rank anything here; ``most-cited work of the period'' is therefore not reported,
because any such ranking would be an artefact of indexing lag.

\clearpage

% ------------------------------------------------------------------ trial watch
\section*{\sffamily\bfseries\color{Deep}Trial registry watch}
\vspace{-1.5mm}{\color{Coral}\rule{\linewidth}{1.1pt}}\vspace{2.5mm}

\footnotesize
Seven registry sweeps ran across August. The accumulated watch list now holds
\textbf{24 trials}, of which \textbf{21} were registered, updated or first seen inside
August. Two new interventional registrations appeared in the month:

\vspace{1.5mm}
\textbf{NCT07111208} (MARKOPOLO/METSGREEN, $n$=180, recruiting), added 4 August, and
\textbf{NCT07797413} --- the \textbf{Savannah Environmental Nutrition Study}, registered
31 August by Georgia Southern University. SENS is a community-based
antioxidant- and nutrient-rich food intervention in residents near the Port of Savannah,
$n$=20, not yet recruiting, starting 1 September 2026. \texttt{analyze\_endpoints} was
run: the \emph{primary} endpoint is \textbf{participant retention} ($\geq$80\% target
through week 12) --- this is a feasibility study, and no efficacy claim should be
attached to it. The secondary endpoints are the interesting part: change from baseline in
\textbf{fractional exhaled nitric oxide} (NIOX Mini, ppb) and \textbf{hydrogen peroxide in
exhaled breath condensate} (RTube collection, Amplex Red fluorescence, $\mu$M), plus
blood pressure, resting heart rate and HbA1c. \textbf{No PM exposure measurement is
registered at all} --- proximity to the port is the entire exposure contrast. That is the
recurring weakness of this registry segment: the intervention is characterised in detail
and the exposure is not characterised at all.

\vspace{1.5mm}
The rest of the August watch list is dominated by \textbf{portable air cleaner and
filtration trials} --- PATHFINDER in heart failure ($n$=400), a multi-centre clean-air RCT
in COPD ($n$=770), NCT05874479 on blood pressure in New York ($n$=440), HEPA in rural
Alaska ($n$=15), ERV plus air cleaner for mould and asthma in low-income housing
($n$=80), and cost-effective filtration with PM monitoring in low-income housing
($n$=55) --- plus three MEDiterranean climate-health resilience arms ($n$=102 each,
outdoor workers, hypertensive 65+, and pregnant women) and INTERLUNG in Bhaktapur
($n$=800). One completed cluster-randomised comparison of HEPA against ionisation/plasma
air cleaners ($n$=180) is the most directly instrument-relevant result to watch for
publication. \textbf{The interventional evidence base for personal and household PM
reduction is thickening in the registry but has not yet reached the literature: four
trial-class records in 556 published papers this month.}

\vspace{4mm}

% -------------------------------------------------------------------- open gaps
\section*{\sffamily\bfseries\color{Deep}Open gaps and read on direction}
\vspace{-1.5mm}{\color{Coral}\rule{\linewidth}{1.1pt}}\vspace{2.5mm}

\begin{keybox}
{\sffamily\bfseries\footnotesize\color{Deep}Gaps in the literature}\\[1.2mm]
\footnotesize
\textbf{(i) Ultrafine measurement still does not exist at network scale.} Twelve
ultrafine records in 556, no low-cost particle-number instrument with a published field
calibration against a CPC, and a scoping review (Rashid et al.) that documents the
epidemiological interest without closing the measurement gap. July's falsifiable
prediction on this point failed; it is re-issued for September.
\textbf{(ii) The calibration reference is itself unvalidated.} Particle-bound water
biases gravimetric mass (Chyzhykov et al.); a full assessment campaign of the gravimetric
method was needed and published (Anon et al.). Every sensor $R^2$ in this corpus is
conditional on a reference nobody has re-audited.
\textbf{(iii) Four trial-class records in 556.} Every causal claim in the month rests on
observational designs with modelled exposure.
\textbf{(iv) Geographic imbalance persists.} 121 China against 16 sub-Saharan Africa and
13 Latin America --- better than July's 44 against 4, still not proportionate to burden,
and Song et al.'s toxicity-weighting perspective puts the highest toxicity-adjusted
exposure precisely in the regions this corpus covers least.
\textbf{(v) Drift and ageing are measured but not modelled.} Chung et al.\ characterise
PurpleAir longevity; Kumpika et al.\ correct sensor ageing; no record in the month puts
an ageing term inside an operational network's correction pipeline.
\end{keybox}

\vspace{2.5mm}

\begin{keybox}
{\sffamily\bfseries\footnotesize\color{Deep}Gaps and defects this pipeline creates}\\[1.2mm]
\footnotesize
\textbf{(i) DUPLICATES SHIPPED --- found by this rollup.} Pooling August returned
\textbf{556 records over 548 unique DOIs}. A full-store audit found \textbf{12 records
published twice}: seven on 2 August (four first shipped 1 August, three first shipped
29 July), two on 3 August, one on 9 August (first shipped 7 August), and the two
\texttt{10.3779/j.issn.1009-3419} records already disclosed in the W35 weekly. In every
case the DOI was \emph{already correctly recorded} in \texttt{state/seen.json} under the
first issue date --- the index was being written and not read, and the duplicate shipped
under a reworded title that defeated the title-signature leg. \textbf{Fixed forward on
1 September}: \texttt{check\_dois.py}, which already gates every \texttt{--post} build,
now hard-fails on \texttt{DUPLICATE-DOI} by scanning the shipped corpus itself rather
than \texttt{seen.json}, so it cannot be defeated by a write-side bug either. Shipped
issues have \emph{not} been retro-edited --- rewriting published issues is worse than
disclosing the error --- so \textbf{every count in this issue is over 556 records
containing 8 within-August duplicates, and subtopic counts are inflated accordingly by
roughly 1.4\%.}
\textbf{(ii) 65 records (12\%) are classification-only.} No abstract was deposited
anywhere reachable at issue time. They are counted but carry no assessed content, and the
absence of criticism on those records is not approval.
\textbf{(iii) No first-author affiliation field}, so no author- or institution-cluster
analysis is possible. Open since 15 August.
\textbf{(iv) Venue names are not normalised} --- Research Square appears as ``Res Sq
(preprint)'' (9), ``Research Square (preprint)'' (9) and ``Research Square'' (3), and
medRxiv as two spellings (5 and 1). Any venue table in this series under-reports split
spellings; the top-ten list above is correct only for venues with a single spelling.
\textbf{(v) OpenAlex has been dead all month} (HTTP 403 paid-plan block). Its coverage is
routed through Consensus, which has no entry-date filter, so that leg contributes dated
backfill rather than windowed harvest.
\end{keybox}

\vspace{3mm}

\footnotesize
\textbf{Read on direction.} July's read was that the field is moving from mass to
composition and number while the instrumentation base stands still. August refines it:
the instrumentation base is moving, but along the \emph{transferability} axis rather than
the \emph{new-quantity} axis. Eight calibration papers, a QA/QC tool, a longevity study
and a distance-uncertainty law all improve how well we can trust a PM$_{2.5}$ mass
network --- and not one of them measures anything a PM$_{2.5}$ mass network does not
already measure. Meanwhile Zhong et al.\ demonstrated inflammatory risk from
photochemical ultrafine particles \emph{at low mass}, and Song et al.\ argued that
toxicity, not mass, should be the control variable. \textbf{The two halves of the field
are optimising different objective functions, and the gap widened in August rather than
narrowing.}
\textbf{Falsifiable predictions for September, both re-issued and sharpened:}
(1) a field calibration of a low-cost particle-number or black-carbon sensor against a
CPC or aethalometer reference appears --- this failed in August and a second miss should
be read as a hardware-cost bottleneck, not a research-attention one; and
(2) at least one paper explicitly incorporates a sensor-\emph{ageing} term into an
operational network correction, rather than characterising ageing and correcting it
offline. If (2) also fails, the practical conclusion for a distributed-sensing
programme is that drift management remains an unclaimed contribution.

\clearpage

% ==================================================================== DIGEST
%% MANDATORY, same contract as the daily and the weekly (formalised 2026-08-02).
%% Full summaries are lifted VERBATIM from issues/digest_*.tex by mkdigest.py.
\section*{\sffamily\bfseries\color{Deep}Paper-level digest --- the month by subtopic}
\vspace{-1.5mm}{\color{Coral}\rule{\linewidth}{1.1pt}}\vspace{2mm}

\footnotesize
Of the 556 pooled records, \textbf{548 carry a full assessed summary} lifted verbatim
from the daily issue that first published it, and \textbf{8 are title-only}. A title-only
entry is \textbf{not an assessment}: it has not been checked against the paper, no
limitation has been identified, and the absence of criticism is not approval. Read
those entries as a pointer to a DOI and nothing more. Note also that the 548 figure
counts \emph{records}, not distinct works --- see defect (i) above.

\vspace{3mm}

\input{digest_body}

\clearpage

% ==================================================================== TABLE
\section*{\sffamily\bfseries\color{Deep}Pooled corpus register --- August 2026}
\vspace{-1.5mm}{\color{Coral}\rule{\linewidth}{1.1pt}}\vspace{2.5mm}

{\fontsize{6.7}{8.4}\selectfont
\setlength{\tabcolsep}{2.4pt}
\begin{longtable}{@{}>{\RaggedRight}p{25mm}>{\RaggedRight}p{41mm}>{\RaggedRight}p{27mm}%
>{\RaggedRight}p{23mm}>{\RaggedRight}p{18mm}>{\RaggedRight}p{15mm}@{}}
\toprule
\textsf{\textbf{First author}} & \textsf{\textbf{Focus}} & \textsf{\textbf{Journal}} &
\textsf{\textbf{Design}} & \textsf{\textbf{Metric}} & \textsf{\textbf{Region}}\\
\midrule
\endfirsthead
\toprule
\textsf{\textbf{First author}} & \textsf{\textbf{Focus}} & \textsf{\textbf{Journal}} &
\textsf{\textbf{Design}} & \textsf{\textbf{Metric}} & \textsf{\textbf{Region}}\\
\midrule
\endhead
\multicolumn{6}{@{}l@{}}{\cellcolor{Mist}\textsf{\textbf{\textcolor{Deep}{Neuro / mental health}}}}\\[0.6mm]
\href{https://doi.org/10.1001/jamanetworkopen.2026.25015}{Dales} & Wildfire-sourced PM$_{2.5}$ and ED visits for migraine and other primary headache & \textit{JAMA Netw Open} & Case-crossover & PM$_{2.5}$ (wildfire) & Canada\\
\href{https://doi.org/10.1093/ageing/afag215}{Glover} & Umbrella meta-review of urban environmental exposures and cognitive health & \textit{Age Ageing} & Umbrella review / meta-analysis & PM$_{2.5}$ & Global\\
\href{https://doi.org/10.1016/j.nsa.2026.107019}{Goossens} & In utero air pollution exposure and pre-/postnatal brain development & \textit{Neurosci Appl} & Narrative review & PM$_{2.5}$ / PM$_{10}$ / PM$_{\mathrm{coarse}}$ & Global\\
\href{https://doi.org/10.1016/j.isci.2026.116617}{Ma} & PM$_{2.5}$-induced neurotoxicity and depression: exposure risks and mechanisms & \textit{iScience} & Narrative review & PM$_{2.5}$ + components & Global\\
\href{https://doi.org/10.3760/cma.j.cn112150-20250804-00752}{Pan} & Maternal PM exposure and offspring neurodevelopment: effects and mechanisms & \textit{Zhonghua Yu Fang Yi Xue Za Zhi} & Narrative review & PM (bulk) & China\\
\addlinespace[0.9mm]
\multicolumn{6}{@{}l@{}}{\cellcolor{Mist}\textsf{\textbf{\textcolor{Teal}{Cardiovascular \& metabolic}}}}\\[0.6mm]
\href{https://doi.org/10.3760/cma.j.cn112338-20251120-00843}{Huang} & PM$_{2.5}$ and dyslipidaemia in rural Chinese children; dietary effect modification & \textit{Zhonghua Liu Xing Bing Xue Za Zhi} & Cross-sectional (multistage) & PM$_{2.5}$ & China\\
\href{https://doi.org/10.32394/pe/218966}{Kazmierska} & Air pollution and cardiovascular disease: review of current data & \textit{Przegl Epidemiol} & Narrative review & PM$_{2.5}$ / PM$_{10}$ & Global\\
\href{https://doi.org/10.1007/s00125-026-06814-2}{Xi} & Air pollution, built environment and temperature vs. prediabetes phenotypes & \textit{Diabetologia} & Prospective cohort & PM$_{2.5}$ & Germany\\
\href{https://doi.org/10.1161/JAHA.125.046917}{Yang} & Short-term PM$_{2.5}$ and hospitalisation/ED use in atrial fibrillation & \textit{J Am Heart Assoc} & Case-crossover & PM$_{2.5}$ (wildfire-influenced) & USA\\
\addlinespace[0.9mm]
\multicolumn{6}{@{}l@{}}{\cellcolor{Mist}\textsf{\textbf{\textcolor{Sky}{Reproductive \& developmental}}}}\\[0.6mm]
\href{https://doi.org/10.1016/j.jes.2025.10.043}{Chen} & Prenatal air pollution, fetal growth trajectory phenotypes and child neurodevelopment & \textit{J Environ Sci (China)} & Prospective cohort & PM$_{2.5}$ / PM$_{10}$ & China\\
\href{https://doi.org/10.1016/j.ecoenv.2026.120566}{Huang} & Joint prenatal bisphenol + air-pollutant mixture exposure and preschool IQ & \textit{Ecotoxicol Environ Saf} & Prospective cohort (mixtures) & PM$_{2.5}$ / PM$_{10}$ & China\\
\href{https://doi.org/10.1007/s40572-026-00552-8}{Peterson} & Ambient air pollution, PAHs and pubertal development: critical review & \textit{Curr Environ Health Rep} & Critical review & PM$_{2.5}$ components & Global\\
\href{https://doi.org/10.1016/j.envres.2026.125313}{Xie} & Long-term PM$_{2.5}$ and timing of pubertal development in boys and girls & \textit{Environ Res} & Prospective cohort & PM$_{2.5}$ & China\\
\addlinespace[0.9mm]
\multicolumn{6}{@{}l@{}}{\cellcolor{Mist}\textsf{\textbf{\textcolor{Sage}{Respiratory \& allergic}}}}\\[0.6mm]
\href{https://doi.org/10.1007/s00484-026-03289-z}{Bayar Muluk} & Climate change and otolaryngology: airway disease impacts and emission reduction & \textit{Int J Biometeorol} & Narrative review & PM$_{2.5}$ & Global\\
\href{https://doi.org/10.4103/mgr.MEDGASRES-D-25-00164}{Liao} & Short-term air pollution, respiratory symptoms and lung function in asthma & \textit{Med Gas Res} & Panel study & PM$_{2.5}$ / PM$_{10}$ & China\\
\href{https://doi.org/10.3390/epidemiologia7040097}{Sanchez de Prada} & Ambient PM and asthma-coded admissions in children 0-3 y, Spain & \textit{Epidemiologia} & Ecological panel & PM$_{2.5}$ / PM$_{10}$ & Spain\\
\href{https://doi.org/10.1016/j.jhazmat.2026.143076}{Yang} & Allergenic vs. toxic risk of airborne seed hairs compared with ambient PM$_{2.5}$ & \textit{J Hazard Mater} & Field measurement & PM$_{2.5}$ + bioaerosol & China\\
\addlinespace[0.9mm]
\multicolumn{6}{@{}l@{}}{\cellcolor{Mist}\textsf{\textbf{\textcolor{Violet}{Mechanistic toxicology}}}}\\[0.6mm]
\href{https://doi.org/10.1016/j.jes.2025.12.074}{Ding} & PM$_{2.5}$ promotes vascular calcification via miR-27b-5p/HES1 signalling & \textit{J Environ Sci (China)} & In vitro + organoid + murine & PM$_{2.5}$ & China\\
\href{https://doi.org/10.1007/s00360-026-01698-5}{Maraschi} & Osmoregulatory disruption and metal bioaccumulation from settleable atmospheric PM & \textit{J Comp Physiol B} & Controlled exposure (ecotox) & Settleable APM & Brazil\\
\href{https://doi.org/10.1016/j.jes.2025.11.018}{Wu} & Macrophage NLRP3 inflammasome in real-ambient PM-induced hepatic glucose dysregulation & \textit{J Environ Sci (China)} & Real-ambient murine + co-culture & Real-ambient PM & China\\
\addlinespace[0.9mm]
\multicolumn{6}{@{}l@{}}{\cellcolor{Mist}\textsf{\textbf{\textcolor{Clay}{Occupational \& indoor}}}}\\[0.6mm]
\href{https://doi.org/10.3389/fmedt.2026.1796828}{Chen} & Nested suction instrument to cut intraoperative surgical-smoke PM$_{2.5}$ & \textit{Front Med Technol} & Bench / device experiment & PM$_{2.5}$ (surgical smoke) & China\\
\href{https://doi.org/10.1093/annweh/wxag055}{Cole} & Determinants of respirable crystalline silica exposure in Australian tunnelling & \textit{Ann Work Expo Health} & Observational exposure study & Respirable dust / RCS & Australia\\
\href{https://doi.org/10.1016/j.jhazmat.2026.143081}{Ding} & Stage-resolved hair Zn-isotope dynamics as a firefighting combustion-exposure tracer & \textit{J Hazard Mater} & Repeated-measures biomarker & Combustion-derived PM & China\\
\href{https://doi.org/10.3390/jox16040119}{Ramsden} & Xenobiotic hazards in aircraft cabin air & \textit{J Xenobiot} & Review + risk economics & Ultrafine particles & Global\\
\addlinespace[0.9mm]
\multicolumn{6}{@{}l@{}}{\cellcolor{Mist}\textsf{\textbf{\textcolor{Amber}{Sensing, forecasting \& instrumentation}}}}\\[0.6mm]
\href{https://doi.org/10.1016/j.jenvman.2026.130515}{Shi} & Adaptive spatiotemporal graph transformer for multi-site multi-step PM$_{2.5}$ forecasting & \textit{J Environ Manage} & Deep-learning model & PM$_{2.5}$ & China / India\\
\href{https://doi.org/10.3390/bios16070380}{Tsuchiya} & Seasonal and indoor-outdoor characteristics of fluorescent aerosol particles in offices & \textit{Biosensors} & Real-time sensor campaign & Bioaerosol / coarse & Japan\\
\addlinespace[0.9mm]
\multicolumn{6}{@{}l@{}}{\cellcolor{Mist}\textsf{\textbf{\textcolor{Amber}{Exposure assessment \& modelling}}}}\\[0.6mm]
\href{https://doi.org/10.1029/2025GH001621}{Solorzano-Araque} & Multiscale exposure modelling: LOTOS-EUROS + mobility surveys + particle toxicity & \textit{GeoHealth} & CTM + mobility integration & PM$_{2.5}$ / PM$_{10}$ & Colombia\\
\addlinespace[0.9mm]
\multicolumn{6}{@{}l@{}}{\cellcolor{Mist}\textsf{\textbf{\textcolor{Coral}{Burden, policy \& mitigation}}}}\\[0.6mm]
\href{https://doi.org/10.1002/hsr2.72886}{Karamzad} & Ischaemic heart disease burden in MENA 1990-2021; PM attributable fraction 28.2\% & \textit{Health Sci Rep} & GBD secondary analysis & Ambient PM & MENA\\
\href{https://doi.org/10.1371/journal.pone.0347989}{Rees} & HEPA filtration to reduce respiratory infection in care homes & \textit{PLoS One} & RCT process evaluation & Filtered PM & UK\\
\href{https://doi.org/10.1016/j.jes.2026.01.009}{Shan} & Policy-driven emission reductions and post-pandemic rebound, high-resolution inventories & \textit{J Environ Sci (China)} & Emission inventory & PM emissions & China\\
\addlinespace[0.9mm]
\multicolumn{6}{@{}l@{}}{\cellcolor{Mist}\textsf{\textbf{\textcolor{Slate}{Other clinical endpoints}}}}\\[0.6mm]
\href{https://doi.org/10.3390/dj14070440}{Lupita} & PM$_{2.5}$, preventive-care capacity and caries/edentulism burden across the EU & \textit{Dent J} & Ecological panel & PM$_{2.5}$ & EU-27\\
\href{https://doi.org/10.1016/j.jdin.2026.06.008}{Park} & Long-term PM exposure and risk of androgenetic alopecia & \textit{JAAD Int} & Nationwide cohort & PM$_{2.5}$ / PM$_{10}$ & South Korea\\
\href{https://doi.org/10.3171/2026.3.PEDS2615}{Sun} & Environmental exposure and area deprivation vs. paediatric hydrocephalus outcomes & \textit{J Neurosurg Pediatr} & Registry cohort + geospatial & PM$_{2.5}$ (satellite-derived) & USA\\
\addlinespace[0.9mm]
\bottomrule
\end{longtable}}


\vspace{4mm}
\begin{keybox}
{\sffamily\bfseries\footnotesize\color{Deep}Provenance and method --- read this before
citing any count above}\\[1.2mm]
{\fontsize{7.2}{8.8}\selectfont
\textbf{This issue keeps the standing rollup contract.} Unlike the July monthly, no API
was re-queried to build it: every one of the 556 records was published in one of the
\textbf{31 daily issues covering 1--31 August 2026}, each of which harvested its own
single-day entry window, and this rollup is pure aggregation over
\texttt{state/corpus/2026-08-*.json} and \texttt{state/metrics.csv}. There is no
backfill batch and no curated sample: subject to the source limitations below, every
proportion here is a proportion of what the daily windows actually recovered.
\\[1.0mm]
\textbf{Sources across the month.} \textbf{PubMed} E-utilities on \texttt{[Date - Entry]},
two axes (health, instrumentation), queried daily. \textbf{Europe PMC} on
\texttt{CREATION\_DATE}; productive early in the month, returning zero on several late-August
days. \textbf{Crossref by-ISSN} across 25 tracked non-PubMed journals, windowed on
\texttt{created} --- this channel was built during August in response to July's stated gap
and is the reason \textit{Atmos Chem Phys} and \textit{Atmos Meas Tech} appear at all.
\textbf{arXiv} relevance query with a local recency screen. \textbf{OpenAlex} returned
HTTP 403 on every call all month and contributed nothing; its coverage was routed through
\textbf{Consensus}, which has no entry-date filter and therefore contributes explicitly
labelled dated backfill rather than windowed harvest. \textbf{Google Scholar} was never a
dependency.
\\[1.0mm]
\textbf{Known limits on every count.} (a) \textbf{Eight within-August duplicate records}
across seven DOIs are included in all totals; see defect (i) in the gaps box. Counts are
over-stated by roughly 1.4\%. (b) \textbf{Sixty-five records (12\%) are
classification-only}, with no abstract available at issue time. (c) Deposit timing at a
few high-volume publishers drives most day-to-day variation, so daily counts are not a
publication-rate signal --- read week totals. (d) Subtopic, design, particle-metric and
geography labels are single-label judgements from title, abstract, keywords and MeSH
terms, not a controlled vocabulary. (e) Effect estimates were transcribed verbatim from
abstracts and only where an interval was reported; 148 qualified, from 63 of 556 records.
(f) Every DOI in every daily issue passed \texttt{check\_dois.py} --- PMID-authority and
resolution checks --- before that issue shipped.
\\[1.0mm]
\textbf{Trial watch} aggregates \texttt{state/trials.json} over seven sweeps; the
ClinicalTrials.gov connector has no reliable last-update filter on some axes, so parts of
the watch are relevance-ranked rather than windowed and the list is a floor.
Content derived from PubMed and Crossref metadata; abstracts remain the copyright of
their respective publishers.\par}
\end{keybox}

\end{document}
